% Pure Math Survey 2.1.0: lecture. Preparation is not completed authorship.
% r5 opening: no abstract or pre-section summary; later section count is editable.
% Current r5 authoring duties: active 2--3 normal TYPESET-line target for explanatory prose.
% Inspect the actual final PDF; record specific long-block reasons, alternatives,
% page/source positions, final PDF hash and reinspection. No TeX line-count games.
% Reconstruct obligations before locating listed steps; save actual fast-route,
% slow-step and condition-change attempts. Exact deep imports still need their
% current hypothesis checks, substitutions and object conversions explained.
% Reusable explanations return to live KB, review and a new snapshot before adoption.
% This scaffold is unfinished; profile/view tasks and substantive evidence remain.
\documentclass[12pt,reqno]{amsart}
\usepackage{math-review}
\input{../generated/conventions.tex}
\title{@@TITLE@@}
\date{}
\begin{document}
\maketitle
@@DRAFT@@
\section{Goal Theorem and Proof Strategy}
\placeholder{State learning goals and the complete core goal theorem early. Give a faithful proof and learning route, and explain which problem each later construction solves. A topic-specific first heading is allowed with the same duties; do not front-load all teaching lemmas.}
\placeholder{Read the selected reviewed overview, routes and node fast passages first; form Question, Main Claim, Mechanisms and Critical Conditions. Use full fixed canonical inputs below. For each borrowed principal place controlling source, version and exact locator in adjacent visible prose; distinguish adapted, local and unverified claims. Paper prior-work and novelty assertions need actual source and slow-detail verification.}
@@PRINCIPALS@@
\placeholder{Choose the teaching order of prerequisites, constructions, calculations, proofs and first applications. State important branches fully near their use; reproduce every promised derivation. Later sections may split or combine; no fixed count.}
\section{Prerequisites and a Motivating Model}
\placeholder{Open only the precise slow subsections required by actual proof promises. Establish consumed important Lemma/Proposition interfaces before their first proof use, separating independent branches. Keep all statement types numbered with labels and no descriptive optional titles; proof-target headings are allowed. Attribute borrowed, modified and composite tools nearby, with actual conversions and every source input; local derivations need no invented citation. Missing mathematics returns to KB, review and a new snapshot.}
@@INGREDIENTS@@
\placeholder{Classify assumed, refreshable and bridge-required inputs relative to this learner, and execute a useful model calculation.}
\section{Construction and Key Estimates}
@@VIEW_TASK@@
\placeholder{Explain the missing output and why the chosen object supplies it before working through promised calculations. Track parameter-independent constants, convergence topology, admissibility and hypotheses at their actual use.}
\section{Proof of the Goal Theorem}
\placeholder{Reconstruct the necessary obligations from the goal and learner before checking the supplied list. Close every promised step; clearly separate exact imported deep results from proofs actually taught here, and explain their present hypothesis checks, substitution and output use.}
\section{Examples and Failure Modes}
\placeholder{Execute the method on a model and explain a genuine boundary. Save actual fast-route reasoning, a reproduced slow calculation and the first failed step after changing a condition; distinguish proof failure from a counterexample. Exercises are optional; no fixed chapter count is imposed.}
\placeholder{Prefer 2--3 normal typeset lines for ordinary explanatory prose; keep continuous arguments together when needed. Rewrite hidden mathematical actions; record page/source positions, specific reasons and alternatives for necessary long blocks, the final PDF hash and reinspection scope. Never use decorative formulas, forced breaks, smaller type or omitted hypotheses.}
\bibliographystyle{plain}
\bibliography{../generated/references}
\end{document}
